\addcontentsline{toc}{section}{\MakeUppercase{Introduction and Research Motivation}}
\section*{\MakeUppercase{Introduction and Research Motivation}}\label{chap:intro}
\glsresetall
\glsunset{etal}
\glsunset{eg}
\glsunset{etc}
\glsunset{ie}
\subsection*{Relevance of the Work}\label{subsec:relevance}
Audio classification plays a crucial role in contemporary \gls{aad} systems by providing {essential insights}\label{rev:ru_min_1a} suitable for various applications~\cite{Kilickaya2025, LoScudo2023, Barusco2025}. These fields are valid not only for industrial environments~\cite{Purohit2019} but also for dynamic urban areas~\cite{Piczak2015, Salamon2014} and even complex natural ecosystems, such as forests~\cite{Bandara2023}. \dots

\dots

\subsection*{Formulation of Problems}\label{sec:form_prob}
Audio-based classification has strong potential in many domains, but recurring issues still obstruct its effective use in real-world data collection and processing. These issues are particularly critical when \dots

\begin{enumerate}[leftmargin=*]
	\item Although handcrafted audio features like \gls{mfcc} and chroma-based representations are widely used, there is still no systematic understanding of \dots
	\item \dots
\end{enumerate}

\subsection*{Aim of the Research}\label{sec:aim}
{This research aims to improve the accuracy, robustness, and reliability of \gls{aad} in noisy environments. It focuses on better \dots}

\subsection*{Research Objectives}\label{sec:r_o}
The overarching goal of this study is to enhance the accuracy, robustness, and dependability of \dots

\begin{enumerate}[leftmargin=*]
    \item To review and integrate prior work on \gls{aad} and \gls{esc}, focusing on feature representations, \gls{dnn} models, and evaluation methods in noisy environments.
    
    \item \dots
\end{enumerate}

\subsection*{Scientific Novelty}\label{subsec:sci_nov}
This research introduces several significant scientific contributions that \dots

\begin{enumerate}[leftmargin=*]
	\item {A key advancement is the full use of the \gls{mimii} dataset across all devices and \gls{snr} levels, unlike prior work restricted to specific devices or \dots}
	\item \dots
\end{enumerate}


\subsection*{Practical Significance}\label{subsec:1_7}
\begin{enumerate}[leftmargin=*]
	\item Applying \gls{dl} to acoustic-based anomaly detection is a promising approach for \dots
	\item \dots
\end{enumerate}

\subsection*{Dissertation Statements} \label{subsec:des_state}
\begin{enumerate}[leftmargin=*]
	\item {The major class imbalance in the \gls{mimii} dataset was addressed using \dots}
	
	\item \dots
\end{enumerate}

\subsection*{Scientific Approvals} \label{sec:scie_app}
{To validate and support the methodology developed in this doctoral research, five scientific contributions have been produced over the past four years. These consist of \dots}

\subsubsection*{Articles indexed in Web of Science and Scopus}
\begin{enumerate}[leftmargin=*]
	\item Ahmad Qurthobi, Rytis Maskeli\={u}nas, and Robertas Dama\v{s}evi\v{c}ius. Detection of mechanical failures in \dots
	\item \dots
\end{enumerate}
\subsubsection*{International conference proceedings}
\begin{enumerate}[leftmargin=*]
	\item Ahmad Qurthobi and Rytis Maskeli\={u}nas. The effect of augmentation and filtration on noisy environment’s acoustic signals to detect abnormalities in \dots
	\item \dots
\end{enumerate}

\subsection*{Dissertation Structure}\label{sec:dis_struc}
This dissertation is structured into five chapters. Chapter~\hyperref[chap:intro]{\MakeUppercase{Introduction and Research Motivation}} presents \dots
